\chapter{Exemples et compl\'ements}
\label{XI}

\section{Espaces projectifs, vari\'et\'es unirationnelles}
\label{XI.1}

\begin{proposition}
\label{XI.1.1}
Soient $k$ un corps alg\'ebriquement clos, $X=\PP_k^r$ l'espace
projectif de dimension $r$ sur~$k$. Alors $X$ est \emph{simplement
connexe,} \ie $\pi_1(X)=0$.
\end{proposition}

Pour $r=0$, c'est trivial. Si $r=1$, il faut montrer que si $X'$ est
un rev\^etement \'etale connexe non vide de~$X=\PP_k^1$, alors
$X'\isomto X$. La formule du genre nous donne ici, si~$g$ et $g'$ sont
les genres de~$X$ et~$X'$:
$$
1-g' = d(1-g)\quoi,
$$
o\`u $d$ est le degr\'e de~$X'$ sur~$X$. Comme $g=0$, on aura donc
$1-g'=d$, ce qui exige $d=1$ puisque $g'\ge 0$, ce qui prouve
$X'\isomto X$. Lorsque $r\ge 2$, on proc\`ede par r\'ecurrence
sur~$r$, en supposant que $\PP^{r'}$ est simplement connexe pour
$r'<r$. Appliquant ceci \`a un hyperplan de~$\PP^r$ et utilisant
(X~\ref{X.2.10}), il en r\'esulte bien que $\PP^r$ est simplement
connexe. Autre d\'emonstration: on aura
$\pi_1(\PP^1\times\cdots\times \PP^1) = \pi_1(\PP^1)\times\cdots
\times \pi_1(\PP^1)$ en vertu de (X~\ref{X.1.7}), donc $(\PP^1)^r$
est simplement connexe puisque $\PP^1$ l'est, donc $\PP^r$ est
simplement connexe en vertu de l'invariance birationnelle du groupe
fondamental (X~\ref{X.3.4}). Cette d\'emonstration montre plus
g\'en\'eralement:

\begin{corollaire}
\label{XI.1.2}
Soit $X$ un sch\'ema propre et normal sur un corps
alg\'ebriquement clos~$k$; si $X$ est une vari\'et\'e
rationnelle, \ie int\`egre et son corps des fonctions est une
extension transcendante pure de~$k$, alors $X$ est simplement connexe.
\end{corollaire}

Ce r\'esultat s'applique en particulier aux vari\'et\'es
grassmanniennes et plus
g\'en\'e\-ra\-le\-ment aux vari\'et\'es~$G/H$, o\`u $G$ est un
groupe lin\'eaire connexe sur~$k$ et $H$ un sous-groupe
alg\'ebrique contenant un sous-groupe de Borel
de~$G$.

Rappelons qu'on appelle vari\'et\'e \emph{unirationnelle sur~$k$}
\index{unirationnelle (vari\'et\'e)|hyperpage}un sch\'ema propre et int\`egre sur~$k$ dont le corps des
fonctions $K$ est contenu dans une extension transcendante pure~$K'$
de~$k$, finie sur~$K$ (\ie ayant m\^eme degr\'e de transcendance
sur~$k$ que~$K$), \ie s'il existe une application rationnelle
dominante $f\colon\PP_k^r\to X$, avec $r=\dim X$. Si $X$ est normale,
on voit donc par les r\'eflexions pr\'ec\'edant X~\ref{X.3.4}
que pour tout rev\^etement \'etale connexe $X'$ de~$X$, de
corps~$L/K$, l'alg\`ebre $L\otimes_K K'$ sur~$K'$ est non
ramifi\'ee sur le mod\`ele~$\PP^r$, donc compl\`etement
d\'ecompos\'ee en vertu de~\ref{X.1.1} ce qui montre que $L$ est
$K$-isomorphe \`a une sous-extension de~$K'/K$. Cela prouve donc,
compte tenu de V~\ref{V.8.2}:

\begin{corollaire}
\label{XI.1.3}
Le groupe fondamental d'une vari\'et\'e unirationnelle normale sur
un corps al\-g\'e\-bri\-quement clos, est fini.
\end{corollaire}

(N.B. On notera que dans la d\'efinition de \og unirationnelle\fg, on
n'avait pas besoin que $K'/K$ soit finie).

\begin{remarques}
\label{XI.1.4}
Bien entendu, les r\'esultats de ce num\'ero sont bien
connus. D'autre part,
J-P.\, Serre a montr\'e~\cite{XI.10} que lorsque $X$ est
une vari\'et\'e projective unirationnelle lisse sur un corps
alg\'ebriquement clos de \emph{caractéristique nulle,} $X$ est
simplement connexe. Sa d\'emonstration est transcendante en ce
qu'elle utilise le th\'eor\`eme de sym\'etrie de Hodge.

\lcrochetbf Ajout\'e en 2003 (MR) Bornons-nous au cas d'un corps alg\'ebriquement clos $k$ de
caract\'eristique $p\geq 0$. En caract\'eristique $p>0$, la d\'efinition de
$k$-vari\'et\'e unirationnelle donn\'ee plus haut est celle de~$k$-vari\'et\'e
faiblement unirationnelle, par opposition \`a celle de~$k$-vari\'et\'e
fortement unirationnelle, o\`u l'on suppose de plus que $K'$ est une
extension s\'eparable de~$K$.

Dans le cas de la dimension $2$, il existe des surfaces projectives et
lisses, faiblement unirationnelles qui ont un groupe fondamental (fini
d'apr\`es 1.3) non trivial et donc qui ne sont pas rationnelles (T.\, Shioda, On unirationality of supersingular surfaces, Math.~Ann. \textbf{225}
(1977), p.\,155--159.). Par contre une surface fortement
unirationnelle est rationnelle: cela r\'esulte du crit\`ere de rationalit\'e de
Castelnuovo, \'etendu \`a la caract\'eristique $p>0$ par O.\, Zariski (\cf J-P.\, Serre, S\'em. Bourbaki \no146, 1957 et {\OE}uvres--Collected
papers, vol.\, 1, p.\, 467).

Sur le corps $\CC$ des nombres complexes, on conna\^it des exemples de vari\'et\'es projectives et
lisses de \hbox{dimension $\geq 3$} qui sont unirationnelles et non
rationnelles (\cf P.\, Deligne, Vari\'et\'es unirationnelles non
rationnelles [d'apr\`es M.\, Artin et D.\, Mumford],  S\'em.\ Bourbaki
\no402, 1971-72, Lecture Notes vol.\, 317). C'est le cas en particulier d'une
hypersurface cubique lisse dans l'espace projectif $\PP^4$.
Certains de ces exemples s'\'etendent en caract\'eristique $>0$ pour donner
des vari\'et\'es fortement unirationnelles non rationnelles (\cf J.P.\, Murre, Reduction of the proof of the non-rationality of a non
singular cubic threefold to a result of Mumford, Compositio Math. \textbf{27}
(1973), p.\,63--82).

Soit $V$ une $k$-vari\'et\'e projective normale, int\`egre. On dit que $V$ est
rationnellement connexe, s'il existe un $k$-sch\'ema de type fini int\`egre
$T$ et un $k$-morphisme $F : T\times\PP^1\to V$ tel que le morphisme :
\begin{equation}\tag*{(1)}
\begin{split}
T \times\PP^1\times\PP^1&\to V\times V\\
(t, u, u')&\mto(F(t,u), F(t,u'))
\end{split}
\end{equation}
soit dominant. En particulier, par deux points rationnels de~$V$,
assez g\'en\'eraux, passe une courbe rationnelle. La terminologie est
justifi\'ee par le fait que si $V$ est rationnellement connexe, on peut
joindre deux points rationnels par une cha\^ine finie de courbes
rationnelles. Si $k$ est de caract\'eristique $p>0$, on renforce la
d\'efinition pr\'ec\'edente en demandant que l'application (1) soit
g\'en\'eriquement lisse. On obtient alors la notion de vari\'et\'e
s\'eparablement rationnellement connexe. Ainsi les vari\'et\'es
unirationnelles sont rationnellement connexes et en caract\'eristique
$p>0$, les vari\'et\'es fortement unirationnelles sont s\'eparablement
rationnellement connexes. J.\, Koll\'ar a montr\'e que les vari\'et\'es
s\'eparablement rationnellement connexes ont un groupe fondamental
alg\'ebrique trivial (\cf O.\, Debarre, Vari\'et\'es rationnellement
connexes [d'apr\`es T.\, Graber, J.\, Harris, J.\, Starr et A.J.\, de~Jong],
S\'em. Bourbaki \no906, 2001-2002).\rcrochetbf
\end{remarques}

\section{Vari\'et\'es ab\'eliennes}
\label{XI.2}
Soient $k$ un corps alg\'ebriquement clos, $A$ une vari\'et\'e
ab\'elienne sur~$k$, \ie un sch\'ema en groupes sur~$k$, propre
sur~$k$, lisse sur~$k$, et connexe, enfin $G$ un sch\'ema en groupes
commutatifs de type fini sur~$k$. D\'esignons par $\Ext(A,G)$ le
groupe des classes d'extensions commutatives de~$A$ par~$G$, par
$\H^1(A,G)$ le groupe des classes de fibr\'es principaux sur~$A$ de
groupe~$G$ (comparer \No \ref{XI.4} plus bas), et consid\'erons
l'homomorphisme canonique
$$
\Ext(A,G)\to \H^1(A,G)\quoi.
$$
Un
raisonnement de Serre \cite[chap.\, VII, th.\, 5]{XI.5} montre que c'est un
homomorphisme injectif, qui a pour image l'ensemble des
\og \'el\'ements primitifs\fg de~$\H^1(A,G)$, \ie des
\'el\'ements $\xi$ pour lesquels on~a:
$$
\pi^*(\xi) = \pr_1^*(\xi) + \pr_2^*(\xi)\quoi,
$$
o\`u $\pr_i$ sont les deux projections de~$A\times A$
sur~$A$, et $\pi\colon A\times A\to A$ la loi de composition de~$A$
(N.B. Serre n'\'enonce son th\'eor\`eme que pour $G$
lin\'eaire et connexe, et bien entendu lisse sur~$k$, mais en
simplifiant la premi\`ere partie de son raisonnement, on voit que
ces restrictions sont inutiles: il suffit de noter que tout morphisme
de~$A$ dans un sch\'ema en groupes $E$ de type fini sur~$k$, qui
transforme unit\'e en unit\'e, est un homomorphisme de groupes, et
d'appliquer ceci aux sections au-dessus de~$A$ d'une extension $E$ de
$A$ par~$G$).

Nous allons appliquer ce r\'esultat au cas o\`u $G$ est un groupe
fini s\'eparable sur~$k$, \ie un groupe fini ordinaire, suppos\'e
commutatif. Utilisant alors $\pi_1(A\times A)\isomto \pi_1(A)\times
\pi_1(A)$ (X~\ref{X.1.7}) et interpr\'etant $\H^1(X,G)$ comme
$\Hom(\pi_1(X),G)$ pour tout sch\'ema alg\'ebrique~$X$, en
particulier pour $X=A$ ou $X=A\times A$, on voit que toute classe de
$\H^1(A,G)$ est primitive, donc on a un isomorphisme
$$
\Ext(A,G)\isomto \H^1(A,G)\quoi,
$$
en d'autres termes \emph{tout revêtement principal de~$A$ de
groupe structural commutatif~$G$, ponctué au-dessus de l'origine
de~$A$, est muni de fa\c con unique d'une structure de groupe
algébrique admettant le point marqué comme origine, et tel que
$A'\to A$ soit un homomorphisme de groupes algébriques.} En
particulier, si $A'$ est connexe, c'est \'egalement une
vari\'et\'e ab\'elienne, isog\`ene \`a~$A$.

D'autre part, comme le foncteur $X\mto \pi_1(X)$ des sch\'emas
alg\'ebriques ponctu\'es $X$ dans les groupes commute au produit
(IX~\ref{IX.1.7}), il transforme un groupe dans la premi\`ere
cat\'egorie en un groupe dans la cat\'egorie des groupes, \ie en
un groupe \emph{commutatif}. Donc \emph{si $A$ est une variété
abélienne, $\pi_1(A)$
est un groupe commutatif.} Donc pour conna\^itre $\pi_1(A)$, il
suffit de conna\^itre le foncteur $G\mto
\H^1(A,G)=\Hom(\pi_1(A),G)$ pour $G$ variant dans les groupes finis
\emph{commutatifs}. Enfin, rappelons que pour tout entier $n>0$,
l'homomorphisme de multiplication par $n$ dans~$A$:
$$
A\lto{n} A
$$
est surjectif, donc \`a noyau fini, \ie c'est une isog\'enie, et
qu'il en r\'esulte que toute isog\'enie $A'\to A$ est quotient
d'une isog\'enie du type pr\'ec\'edent. De ceci, et de
raisonnements standard (\cf par exemple~\cite{XI.6}) on tire:

\begin{theoreme}[Serre-Lang]
\label{XI.2.1}
Soit $A$ une vari\'et\'e ab\'elienne sur un corps
alg\'ebriquement clos~$k$, et pour tout entier $n>0$ consid\'erons
le groupe fini ordinaire $K_n$ sous-jacent au noyau ${}_nA$ de la
multiplication par $n$ dans~$A$, enfin posons pour tout nombre
premier~$\ell$:
$$
T_\ell(A)=\varprojlim_r K_{\ell^r}
$$
et
$$
T_{.}(A)=\underset{\ell}{\prod}T_\ell(A)=\varprojlim_n K_n
$$
(o\`u pour $m$ multiple de~$n$, $m=ns$, on envoie $K_m$ dans $K_n$
par la multiplication par~$s$). Alors le groupe $\pi_1(A)$ est
canoniquement isomorphe \`a $T_{.}(A)$, donc pour tout nombre
premier $\ell$, la composante $\ell$-primaire de~$\pi_1(A)$ est
canoniquement isomorphe \`a $T_\ell(A)$.
\end{theoreme}

On notera que ces isomorphismes sont fonctoriels pour $A$ variable. Le
module $T_\ell(A)$ est appel\'e le \emph{module $\ell$-adique de
Tate} de la vari\'et\'e ab\'elienne~$A$. C'est un foncteur
additif en~$A$, en particulier il donne lieu \`a une
repr\'esentation de l'anneau $\Hom(A,A)$ des endomorphismes de~$A$
dans $T_\ell(A)$, appel\'ee \emph{représentation $\ell$-adique
de Weil}, et qui joue un r\^ole important dans la th\'eorie des
vari\'et\'es ab\'eliennes (\cf par exemple \hbox{\cite[chap~VII]{XI.4}}). Le
th\'eor\`eme~\ref{XI.2.1} en donne une interpr\'etation en
termes de la repr\'esentation naturelle dans le \emph{groupe
d'homologie $\ell$-adique} de~$A$, $\H_1(A,\ZZ_\ell)=\pi_1(A)_\ell$,
ce
qui est \'evidemment plus satisfaisant a priori, du point de vue
notamment de la formule de Lefschetz, \cite[Chap~V]{XI.4}. Rappelons ici les
r\'esultats de Weil sur la structure de~$T_\ell(A)$:
\begin{enumerate}
\item[a)] Si $n$ est premier \`a $\car(k)$, alors $K_n$ est un module
libre de rang \'egal \`a $2\dim A$ sur~$\ZZ/n\ZZ$, donc si $\ell$
est un nombre premier $\neq \car(k)$, $T_\ell(A)$ est un module libre
de rang \'egal \`a $2\dim A$ sur l'anneau $\ZZ_\ell$ des entiers
$\ell$-adiques;
\item[b)] Si $n$ est une puissance de~$\car(k)=p$, alors $K_n$ est un
module libre de rang $\nu\leq\dim A$ sur~$\ZZ/n\ZZ$, $\nu$
ind\'ependant de~$n$, donc $T_p(A)$ est un module libre de rang
$\nu\leq \dim A$ sur l'anneau $\ZZ_p$ des entiers $p$-adiques.
\end{enumerate}

Cela montre que dans la th\'eorie du groupe fondamental
d\'evelopp\'ee ici, le groupe fondamental d'une vari\'et\'e
ab\'elienne variable ne varie pas de fa\c con r\'eguli\`ere
avec le param\`etre, sa composante $p$-primaire pouvant diminuer
brusquement pour des valeurs du param\`etre $t$ correspondant \`a
une caract\'eristique r\'esiduelle~$p$; le cas le mieux connu de
ce ph\'enom\`ene est celui des courbes elliptiques. On notera
cependant que, quel que soit $n$ (premier ou non \`a la
caract\'eristique) le vrai noyau ${}_nA$ dans $A$ pour la
multiplication par $n$ est un sch\'ema en groupe fini sur~$k$ de
degr\'e $n^{2g}$, o\`u $g=\dim A$, qui sera non s\'eparable sur~$k$ si $n$ est un multiple de~$p=\car(k)$. D'ailleurs, lorsque $A$
varie dans une famille de vari\'et\'es ab\'eliennes, \ie si on
a un sch\'ema ab\'elien $A$ sur un sch\'ema de base $S$, on
montre plus g\'en\'eralement que ${}_n A$ est un sch\'ema en
groupes fini et plat sur~$S$, de degr\'e $n^{2g}$ sur~$S$,
c'est-\`a-dire \`a condition de tenir compte des parties
infinit\'esimales des noyaux~${}_n A$, ils se comportent de fa\c con r\'eguli\`ere quel que soit~$n$. Cela sugg\`ere que le
\og vrai\fg groupe fondamental d'une vari\'et\'e ab\'elienne $A$
est le pro-groupe alg\'ebrique (limite projective formelle de
groupes finis sur~$k$) $\varprojlim_n{{}_nA}$, o\`u
par \og vrai groupe fondamental\fg d'un sch\'ema alg\'ebrique~$X$, il
faut entendre: le pro-groupe qui classifie les rev\^etements
principaux de~$X$ de groupe structural un groupe fini quelconque $G$
sur~$k$ (pas n\'ecessairement s\'eparable sur~$k$). De cette
fa\c con par exemple, on r\'ecup\`ere par les
repr\'esentations de~$\Hom(A,A)$ dans la composante $p$-primaire du
vrai groupe fondamental de~$A$, le polyn\^ome caract\'eristique de
Weil d\'efini par ce dernier \`a l'aide des $\ell\neq p$, de fa\c con plus naturelle que la construction de Serre~\cite{XI.8}.

\section{C\^ones projetants, exemple de Zariski}
\label{XI.3}

Soit toujours $k$ un corps alg\'ebriquement clos pour simplifier, et
soit $V$ un $k$-sch\'ema projectif connexe, sous-sch\'ema
ferm\'e de~$\PP^r_k$, qu'on pourra si on veut supposer non
singulier. Soient $Y=\hat C$ le c\^one projetant projectif de~$V$,
$y_0$ son sommet, $X=\hat C_V$ la fermeture projective habituelle du
fibr\'e vectoriel $C_V=\VV(\cal{O}_V(1))$ associ\'e \`a
$\cal{O}_V(1)$, enfin
$$
f\colon X\to Y
$$
le morphisme canonique, contractant la section nulle $X_0$ de~$C_V$
sur~$X$ en un point (EGA~II~8.6.4). Comme $X$ est un fibr\'e
localement trivial sur~$V$, de fibres $\PP^1$ donc de fibres
simplement connexes, le morphisme $p\colon X\to V$ induit en vertu
de~(X~\ref{X.1.4})
un isomorphisme:
$$
\pi_1(X)\isomto\pi_1(V)\quoi.
$$
Comme $p$ induit un isomorphisme $X_0\to V$, on en conclut qu'\emph{un
revêtement étale de~$X$ est complètement décomposé
si et seulement si sa restriction à $X_0$ l'est}. Or pour tout
rev\^etement \'etale $Y'$ de~$Y$, l'image inverse $X'=X\times_Y
Y'$ est un rev\^etement \'etale de~$X$ compl\`etement
d\'ecompos\'e sur la fibre~$X_0$, donc trivial. Comme
l'homomorphisme $\pi_1(X)\to\pi_1(Y)$ est surjectif (IX~\ref{IX.3.4}),
on en conclut que
$$
\pi_1(Y)=(e)
$$
en d'autres termes \emph{tout cône projetant projectif est
simplement connexe.} (N.B. en ca\-rac\-t\'e\-ris\-tique~$0$, le
m\^eme r\'esultat sera valable en prenant pour $Y$ le c\^one
projetant affine).

Supposons maintenant $V$ r\'eguli\`ere \ie lisse sur~$k$; alors
$X$ est r\'eguli\`ere, et pour une immersion projective convenable
de~$V$, on trouve alors un c\^one projetant $Y$ \emph{normal}. Si
$V$ n'est pas simplement connexe, donc $X$ non simplement connexe,
soit $X'$ un rev\^etement \'etale connexe non trivial
de~$X$. Comme les fibres de~$X$ en les points $y\in Y$ distincts de
$y_0$ sont r\'eduites \`a un point, on voit que la restriction de
$X'$ \`a ses fibres (en particulier \`a la
fibre
g\'en\'erique) est triviale; cependant $X$ ne provient pas par
image inverse d'un rev\^etement \'etale de~$Y$, puisque $Y$ est
simplement connexe et que $X'$ serait compl\`etement
d\'ecompos\'e. Cela montre que (X~\ref{X.1.3} et~\ref{X.1.4})
deviennent faux si on remplace l'hypoth\`ese que $f$ est
s\'eparable par celle plus faible que ses fibres sont des
sch\'emas alg\'ebriques s\'eparables (ou m\^eme lisses) sur
les~$\kres(s)$. On notera de m\^eme que les groupes fondamentaux des
fibres g\'eom\'etriques $\overline X_y$ des $y\neq y_0$ sont
\'evidemment r\'eduits \`a $(e)$ puisque ces fibres sont
r\'eduites \`a un point, alors que $\pi_1(X_0)\neq e$, donc le
th\'eor\`eme de semi-continuit\'e (X~\ref{X.2.4}) est
\'egalement en d\'efaut pour~$f$.

Indiquons enfin l'exemple, signal\'e par Zariski, mettant en
d\'efaut ces m\^emes th\'eor\`emes, lorsqu'on y remplace
l'hypoth\`ese que $f$ est s\'eparable par celle que $f$ est
plate. Soit $f\colon X\to Y$ un morphisme d'une surface non
singuli\`ere projective dans la droite rationnelle $Y=\PP^1$, tel
que
$K=k(X)$
soit une extension \og r\'eguli\`ere\fg \ie primaire et
s\'eparable de~$k(f)$, (\ie la fibre g\'en\'erique
g\'eom\'etrique est connexe et s\'eparable), et telle que le
diviseur $(f)=X_0-X_\infty$ soit un multiple
\nieme
d'un diviseur (o\`u $n$ est un entier premier \`a la
caract\'eristique). Il est possible de construire de tels exemples
en toute caract\'eristique. Soit $X'$ le normalis\'e de~$X$ dans
$K(f^{1/n})$, o\`u $K=k(X)$ est le corps des fonctions de~$X$. Il
r\'esulte de l'hypoth\`ese sur~$(f)$ que $X'$ est \'etale
sur~$X$. Soit $Y'$ le normalis\'e de~$Y$ dans $k(t)(t^{1/n})$, il
est ramifi\'e sur~$Y$ en les points $t=0$ et $t=\infty$ exactement,
et la restriction $X'|f^{-1}(U)$ est isomorphe \`a l'image inverse
de~$Y'|U$. En particulier, la restriction de~$X'$ \`a la fibre
g\'en\'erique \emph{géométrique} de~$X$ sur~$Y$ se
d\'ecompose compl\`etement. Cependant, $X'$ n'est pas isomorphe
\`a l'image inverse d'un rev\^etement \'etale de~$Y$, car on
voit tout de suite que ce dernier serait n\'ecessairement~$Y'$, ce
qui est absurde puisque $Y'$ est ramifi\'e sur~$Y$\footnote{On peut
remarquer, du point de vue de la \og topologie \'etale\fg (SGA~4~VII),
que dans cet exemple $\R^1(f_{\et})_*(\ZZ/n\ZZ)$ est \og non
s\'epar\'e\fg sur~$S$.}.

Voici (d'apr\`es Serre) une fa\c con simple de r\'ealiser les
conditions de cet exemple, en s'inspirant de \cite[\No 20]{XI.5}: on prend
pour $n$ un nombre premier $\geq 5$, distinct de la
caract\'eristique, et on fait op\'erer $G=\ZZ/n\ZZ$ dans
$\underline{k}^4${\renewcommand{\thefootnote}{*}\addtocounter{footnote}{-1}\footnote{\lcrochetbf Ajout\'e en 2003: $\underline{k}^4$ d\'esigne l'espace affine de dimension $4$ sur $k$.\rcrochetbf}}
en multipliant les coordonn\'ees par quatre
caract\`eres distincts de~$G$ (ce qui est possible puisque $n\geq
5$). Alors $G$ op\`ere sur l'espace projectif $\PP^3_k$, et les
seuls points fixes sous $G$ sont les quatre points correspondants aux
axes de coordonn\'ees. La surface $X'$ d'\'equation
$x^n+y^n+z^n+t^n=0$
est lisse sur
$k$ (crit\`ere jacobien), et ne contient aucun des points fixes,
donc $G$ \'etant d'ordre premier, op\`ere sur~$X$ \og sans points
fixes\fg \ie $X$ est un rev\^etement principal de~$X=X'/G$ de groupe
$G$. Soit $g=x/y$ dans $k(X')=K'$, c'est un g\'en\'erateur
kumm\'erien de~$K'$ sur~$K=k(X)$ si les caract\`eres choisis
\'etaient~$\chi^i$, $i=0,1,2,3$, avec $\chi$ un caract\`ere
primitif, soit $f$ sa puissance \nieme, qui est un
\'el\'ement de~$K$. On voit tout de suite que $K'$ est une
extension r\'eguli\`ere de~$k(g)$, ce qui r\'esulte du fait que
la courbe plane d'\'equation homog\`ene en $U,T,Z$:
$T^n+Z^n+(1+g^n)U^n=0$ est lisse sur~$k(g)$ (crit\`ere jacobien), et
qu'on sait que toute courbe plane est connexe. D'autre part on a
$k(f)=K\cap k(g)$, puisque le deuxi\`eme membre est une extension de
$k(f)$ contenue dans l'extension $k(g)$ de degr\'e premier, et
distincts de~$k(g)$ (puisque $g\not\in K$). Cela implique que $K$ est
une extension r\'eguli\`ere de~$k(f)$. Enfin le diviseur de~$f$
sur~$X$ est un multiple \nieme d'un diviseur, car son image
inverse sur~$X'$ est le diviseur de~$g^n$, donc un multiple
\nieme, et on peut redescendre parce que $X'$ est \'etale sur~$X$. On aurait fini si l'application rationnelle $f\colon X\to \PP^1$
\'etait un morphisme, c'est-\`a-dire si les diviseurs des
z\'eros et des p\^oles de~$f$ ne se rencontraient point. En fait,
on v\'erifie ais\'ement (en regardant encore sur~$X'$) que les
deux diviseurs en question sont les produits par $n$ de deux courbes
lisses sur~$k$, se coupant transversalement en un point~$a$. Rempla\c cant maintenant $X$ par le sch\'ema obtenu en faisant \'eclater~$a$, les conditions pr\'ec\'edentes ($\divisor(f)$ divisible par
$n$, et $k(X_1)=k(X)$ extension r\'eguli\`ere de~$k(X)$) restent
v\'erifi\'ees, mais de plus $f$ est un \emph{morphisme} $X_1\to
\PP^1$, donc on est sous les conditions voulues.

\section{La suite exacte de cohomologie}
\label{XI.4}

Soit $S$ un pr\'esch\'ema, de sorte que la cat\'egorie
$\Sch_{/S}$ des pr\'esch\'emas sur~$X$ est d\'etermin\'ee,
donc aussi la notion de groupe dans icelle, qu'on appellera aussi
\emph{préschéma en groupes sur~$S$}, ou simplement
\emph{$S$-groupe}. Pour simplifier l'exposition et fixer les
id\'ees, nous nous bornerons le plus souvent par la suite \`a des
groupes qui sont \emph{affines} et \emph{plats} sur~$S$\footnote{En
fait, pour ce qui va suivre, l'hypoth\`ese quasi-affine au lieu
d'affine suffirait, \cf note de bas de page~\eqref{note-296}, p.\,\pageref{note-296}.}, ce qui
suffira pour les applications que nous avons en vue. (Bien entendu, on
rencontre de nombreux cas o\`u ni l'une ni l'autre hypoth\`ese
n'est v\'erifi\'ee). Soit
$G$ un tel $S$-groupe, et soit $P$ un pr\'esch\'ema sur~$S$ sur
lequel $G$ op\`ere \`a droite, ce qui implique en particulier un
morphisme
$$
\pi:P\times_SG\to P
$$
dans la cat\'egorie $\Sch_{/S}$, satisfaisant les axiomes bien
connus. On dit que $P$ est \emph{formellement principal homogène
sous} $G$ si le morphisme
$$
P\times_S G\to P\times_S P
$$
de composantes $\pr_1$ et $\pi$ est un isomorphisme; il revient au
m\^eme de dire que pour tout objet $S'$ de~$\Sch_{/S'}$,
$P(S')=\Hom_S(S',P)$ consid\'er\'e comme ensemble \`a groupes
d'op\'erateurs $G(S')=\Hom_S(S',G)$, est vide ou principal
homog\`ene (\ie vide ou isomorphe \`a $G(S')$ sur lequel le
groupe $G(S')$ op\`ere par translations \`a droite). On dit que
$P$ est \emph{trivial} si $P$ est isomorphe \`a $G$, sur lequel $G$
op\`ere par translations \`a droite, ou ce qui revient au
m\^eme, si chacun des ensembles \`a op\'erateurs $P(S')$ sous
$G(S')$ est trivial. On v\'erifie, par exemple par le
proc\'ed\'e brevet\'e de passage au cas ensembliste, que $P$
\emph{est trivial si et seulement
s'il
est formellement principal homogène, et admet une section sur~$S$}
(ce dernier fait s'\'enon\c cant en termes cat\'egoriques en
disant que $P$ a une section sur l'objet final $e=S$ de~$\Sch_{/S}$,
\ie qu'il existe un morphisme de~$e$ dans $P$). Pour d\'efinir la
notion de fibr\'e principal homog\`ene $P$ sous $G$, plus forte
que celle de fibr\'e formellement principal homog\`ene, il faut
pr\'eciser d'abord dans $\Sch_{/S}$ un ensemble de morphismes qui
seront utilis\'es pour la \og descente\fg, et joueront le r\^ole de
\og morphismes de localisation\fg pour \og trivialiser\fg des fibr\'es. Le
choix le plus ad\'equat varie suivant le contexte, aucun ne
contenant tous les autres\footnote{\label{note-293}Voir \`a ce sujet
SGA~3~IV, notamment~\S 4.}. Ici, il sera commode d'adopter la
d\'efinition suivante:

\begin{definition}
\label{XI.4.1}
Soit $G$ un $S$-groupe. On appelle \emph{fibré principal
homogène} (\`a droite) sous~$G$, un $S$-pr\'esch\'ema $P$
\`a~$S$-groupe \`a droite~$G$, tel qu'il existe un recouvrement de~$S$ par des ouverts~$U_i$, et pour tout $i$ un morphisme de changement
de base $S'_i\to U_i$ fid\`element plat et quasi-compact, tel que
$P'=P\times_S S'$ soit un pr\'esch\'ema \`a op\'erateurs
trivial sous $G'=G\times_S S'$, o\`u $S'$ est le
$S$-pr\'esch\'ema somme disjointe des~$S'_i$.
\end{definition}

(On
notera que le foncteur changement de base $X\mto X'=X\times_S
S'$
\'etant exact \`a gauche, transforme groupes en groupes, objets
\`a groupe d'op\'erateurs en objets \`a groupes
d'op\'erateurs). Notons que \ref{XI.4.1} est \emph{stable par
changement de base}. Notons aussi:

\begin{proposition}
\label{XI.4.2}
Soient $G$ un $S$-groupe, plat et quasi-compact sur~$S$, $P$ un
$S$-pr\'esch\'ema o\`u $G$ op\`ere \`a droite. Conditions
\'equivalentes:
\begin{enumerate}
\item[(i)] $P$ est un fibr\'e principal homog\`ene
sous~$G$;
\item[(ii)] $P$ est formellement principal homog\`ene sous~$G$,
\index{formellement principal homog\`ene (pr\'esch\'ema)|hyperpage}et le morphisme structural $P\to S$ est fid\`element plat et
quasi-compact.
\end{enumerate}
\end{proposition}

Si $P$ est principal homog\`ene sous~$G$, alors avec les notations
de~\ref{XI.4.1} $P'$ est fid\`element plat et quasi-compact sur~$S'$
(puisque $G'$ l'est, et $P'$ lui est $S'$-isomorphe), donc $P$ a les
m\^emes propri\'et\'es au-dessus de~$S$, (pour \og surjectif\fg et
\og quasi-compact\fg, \cf VIII~\ref{VIII.3.1}, pour \og plat\fg c'est un
oubli dans les sorites de l'expos\'e~\ref{VIII}). Inversement, si
(ii) est v\'erifi\'e, prenons le changement de base $S'=P$, qui
est bien fid\`element plat et quasi-compact sur~$S$; alors $P'$ sera
formellement principal homog\`ene sur~$S'$ puisque $P$ l'est sur~$S$
et que le foncteur changement de base est exact \`a gauche, d'autre
part $P'$ a une section sur~$S'$,
\`a savoir
la section diagonale, donc
c'est un fibr\'e principal trivial, ce qui ach\`eve la
d\'emonstration.

\begin{corollaire}
\label{XI.4.3}
Si $G$ est affine et plat sur~$S$, tout fibr\'e principal
homog\`ene $P$ sous $G$ est affine et plat sur~$S$.
\end{corollaire}

En effet, il le devient par extension fid\`element plate et
quasi-compacte de la base, et on applique (VIII~\ref{VIII.5.6}).

L'utilit\'e de la d\'efinition~\ref{XI.4.1} pour des $S$-groupes
\emph{plats} et \emph{affines} sur~$S$ tient \`a
(VIII~\ref{VIII.2.1}), \ie au fait que les morphismes $S'\to S$
envisag\'es dans~\ref{XI.4.1} sont des morphismes de descente
effective pour la cat\'egorie des pr\'esch\'emas affines sur
d'autres. Gr\^ace \`a ce fait, la v\'erification des faits
esquiss\'es ci-dessous se fait de fa\c con essentiellement
\og cat\'egorique\fg\footnote{\Cf \loccit dans note de bas de
page~\eqref{note-293}, p.\,\pageref{note-293}.}. Soit $E$ un $S$-pr\'esch\'ema sur lequel le
$S$-groupe $G$ op\`ere \`a gauche, et soit $P$ un fibr\'e
principal homog\`ene (\`a droite) sous~$G$, nous voulons
d\'efinir un fibr\'e associ\'e $E^{(P)}$, \og localement\fg
isomorphe \`a~$E$. Pour ceci, faisons op\'erer
\`a droite $G$ dans $P\times_S E$ suivant la loi $(x,y)\mto
(xg,g^{-1}y)$, qui d\'ecrit de telles op\'erations dans le
contexte ensembliste, et s'\'etend aux cat\'egories par le
proc\'ed\'e brevet\'e. On posera, sous-r\'eserve d'existence:
$$
E^{(P)}=(P\times_S E)/G
$$
moyennant quoi on constate que $P\times_S E$ sera un pr\'esch\'ema
au-dessus de~$T=E^{(P)}$, \`a groupes d'op\'erateurs \`a droite
$G_T=G\times_S T$; pour \^etre \`a l'aise, on aimerait que de plus
$P\times_S E$ soit un fibr\'e principal homog\`ene sur~$T$ de
groupe $G_T$. Pour v\'erifier l'existence de~$E^{(P)}$ et la
propri\'et\'e pr\'ec\'edente, reprenons le $S'$ de la
d\'efinition~\ref{XI.4.1} et regardons la situation image inverse
sur~$S'$ de la situation initiale. Du fait que $P'$ est
trivial,
\ie isomorphe \`a $G'_d$, on voit tout de suite que $E^{\prime(P')}$
existe, et a la propri\'et\'e voulue d'exactitude. En fait,
$E'\times_{S'} P'$ est $G'$-isomorphe au produit $E'\times_{S'} G'$,
donc $E^{\prime(P')}$ est isomorphe \`a~$E'$. De plus, la formation du
\og fibr\'e associ\'e\fg dans le cas d'un espace \`a op\'erateurs
trivial commute \`a toute extension de la base, et prenant en
occurrence les diverses extensions de la base
$\xymatrix@C=.5cm{{S''}\ar@<2pt>[r]\ar@<-2pt>[r]&{S'}}$
et
$\xymatrix@C=.5cm{{S'''}\ar@<4pt>[r]\ar@<0pt>[r]\ar@<-4pt>[r]&{S'}}$, o\`u
$S''$ et $S'''$ sont les produits fibr\'es double et triple de~$S'$
sur~$S$, on constate que $E^{\prime(P')}$ \emph{est muni d'une donnée
de descente relativement au morphisme $S'\to S$, et que $E^{(P)}$
existe avec les propriétés requises si et seulement si cette
donnée de descente est effective}; bien entendu $E^{(P)}$ n'est
autre alors que l'objet descendu. (Utiliser le fait que $S'\to S$ est
un morphisme de descente dans la cat\'egorie des
$S$-pr\'esch\'emas, \cf VIII~\ref{VIII.5.2}). Il s'ensuit~que
\emph{le fibré associé existe si $E$ est affine sur~$S$}. Nous
appliquerons cette construction au cas o\`u on~a un homomorphisme de~$S$-groupes $G\to H$, et qu'on prend pour $E$ le $S$-pr\'esch\'ema~$H$ muni des op\'erations de~$G$ sur~$H$ \`a gauche r\'esultant
du morphisme donn\'e; comme $H$~op\`ere \`a droite sur
lui-m\^eme de fa\c con \`a commuter aux op\'erations de~$G$
sur~$H$, et que (sous r\'eserve d'existence au-dessus de~$S$) la
formation du fibr\'e associ\'e commute \`a l'extension de la
base, on constate ais\'ement que $H$ va op\'erer \`a droite sur~$P^{(H)}$, qui est d\`es lors un fibr\'e principal homog\`ene
sous $H$ au sens de~\ref{XI.4.1}, et
de fa\c con pr\'ecise est trivialis\'ee par le m\^eme
morphisme $S'\to S$ que~$P$. En particulier, \emph{à tout
fibré principal homogène $P$ sous $G$ et tout homomorphisme de~$S$-groupes $G\to H$, avec $H$ affine sur~$S$, est associé un
fibré principal homogène de groupe~$H$}, de fa\c con
fonctorielle en $(G\to H)$, et compatible avec les changements de base
quelconques $T\to S$.

\begin{definition}
\label{XI.4.4}
Soit $G$ un $S$-pr\'esch\'ema. On note $\H^0(S,G)$ l'ensemble des
sections de~$G$ sur~$S$, qu'on consid\'erera comme un groupe lorsque
$G$ est un $S$-groupe. Dans ce cas, on note $\H^1(S,G)$ l'ensemble des
classes, \`a isomorphisme pr\`es, de fibr\'es principaux
homog\`enes sur~$S$ de groupe $S$, en consid\'erant $\H^1(S,G)$
comme muni du \og point marqu\'e\fg qui correspond aux fibr\'es
triviaux\footnote{\label{note-296}Cette notation n'est coh\'erente
vis-\`a-vis
des notations cohomologiques g\'en\'erales (SGA~4~V)
que lorsqu'on dispose de crit\`eres d'effectivit\'e de descente,
qui ne sont gu\`ere assur\'es que si $G$ est affine (ou seulement
quasi-affine, \cf (VIII~\ref{VIII.7.9})).}.
\end{definition}

Ainsi, $\H^0(S,G)$ est un foncteur en le $S$-pr\'esch\'ema~$G$,
\`a valeurs dans la cat\'egorie des ensembles. Ce foncteur est
exact \`a gauche, a fortiori commute aux produits finis, ce qui
implique en effet qu'il transforme groupes en groupes, groupes
commutatifs en groupes commutatifs. De fa\c con analogue,
$\H^1(S,G)$ est un foncteur en le $S$-groupe \emph{affine}~$G$, \`a
valeurs dans la cat\'egorie des ensembles gr\^ace \`a la
formation des fibr\'es associ\'es; on constate facilement que ce
foncteur commute aux produits finis. En particulier il transforme les
groupes dans la cat\'egorie des $S$-groupes affines, \ie les
$S$-groupes \emph{affines commutatifs}, en des groupes de la seconde,
et m\^eme en des groupes commutatifs (puisque les groupes de la
premi\`ere cat\'egorie sont commutatifs). Ainsi, \emph{si $G$ est
un $S$-groupe affine commutatif, $\H^1(S,G)$ est un groupe
commutatif}, et un homomorphisme $G\to H$ de~$S$-groupes affines
commutatifs donne naissance \`a un homomorphisme de groupes
$\H^1(S,G)\to\H^1(S,H)$.

Pour simplifier, nous nous bornons pour la suite \`a la
consid\'eration de~$S$-groupes \emph{affines et commutatifs}. Soit
$$
0\to G'\lto{u} G\lto{v}G''\to 0
$$
une suite de morphismes de tels groupes, \emph{nous dirons que cette
suite est exacte si $vu=0$} (ce qui permet de consid\'erer $G$
comme un pr\'esch\'ema sur~$G''$, \`a groupes d'op\'erateurs
\`a droite~$G'_{G''}$) \emph{et si $G$ est un fibré principal
homogène
sur~$G''$ de groupe $G'_{G''}=G'\times_S G''$}. Cela implique en
particulier que $u\colon G'\to G$ est un noyau de~$v$, et a fortiori
cela implique l'exactitude de la suite
$$0\to\H^0(X,G')\to\H^0(X,G)\to\H^0(X,G'').$$ Cela implique de plus la
possibilit\'e de d\'efinir une application
$$
\partial\colon\H^0(X,G'')\to\H^1(X,G')\quoi,
$$
en associant \`a toute section de~$G''$ sur~$S$, \ie \`a tout
$S$-morphisme $f\colon S\to G''$, le fibr\'e principal homog\`ene
$P_f$ de groupe $G'\simeq f^*(G'_{G''})$ sur~$S$, image inverse du
fibr\'e principal homog\`ene $G$ sur~$G''$. Du point de vue
$S$-pr\'esch\'emas, ce n'est donc autre que l'image inverse par
$v\colon G\to G''$ du sous-pr\'esch\'ema image de~$S$ par
l'immersion~$f$, et les op\'erations de~$G'$ sur~$P_f$ sont induites
par les op\'erations \`a droite de~$G'$
sur~$G$.
Nous laissons \'egalement au lecteur la v\'erification de la
proposition suivante, qui ne pr\'esente pas de difficult\'es
autres que de r\'edaction:

\begin{proposition}
\label{XI.4.5}
L'application $\partial\colon \H^0(X,G'')\to\H^1(X,G')$ est un
homomorphisme de groupes. La suite d'homomorphismes suivante est
exacte:
\begin{multline*}
0\to\H^0(X,G')\lto{u^0}\H^0(X,G)\lto{v^0}\H^0(X,G')\lto{\partial}\H^1(X,G'')\\
\lto{u^1}\H^1(X,G) \lto{v^1}\H^1(X,G'')
\end{multline*}
(o\`u les homomorphismes autres que $\partial$ proviennent de la loi
fonctorielle de~$\H^0$ \resp~$\H^1$).
\end{proposition}

\begin{remarques}
\label{XI.4.6}
Le point de vue expos\'e ici pour l'\'etude des fibr\'es
principaux homog\`enes est visiblement inspir\'e de Serre~\cite{XI.7}, que
le lecteur aura tout int\'er\^et \`a consulter. Lorsqu'on veut
un formalisme qui s'applique \'egalement \`a des $S$-groupes
structuraux qui sont quasi-projectifs sur~$S$ (de fa\c con \`a
englober les sch\'emas ab\'eliens projectifs en particulier), on a
int\'er\^et \`a modifier \ref{XI.4.1} en y demandant que $S'$
soit somme de pr\'esch\'emas $S'_i$ qui sont finis et localement
libres sur des ouverts $S_i$ de~$S$ recouvrant $S$. Les
d\'eveloppements pr\'ec\'edents sont alors valables, y inclus
notamment \ref{XI.4.5}, en rempla\c cant partout l'hypoth\`ese
affine par l'hypoth\`ese quasi-projective, et en interpr\'etant de
fa\c con correspondante la d\'efinition donn\'ee plus faut d'une
suite exacte de~$S$-groupes. Il suffit en effet de remplacer la
r\'ef\'erence \`a (VIII~\ref{VIII.2.1})
par (VIII~\ref{VIII.7.7}): les morphisme utilis\'es $S'\to S$ sont
des morphismes de descente effective pour la cat\'egorie fibr\'ee
des pr\'esch\'emas quasi-projectifs sur d'autres. On fera
attention cependant que cette deuxi\`eme notion de fibr\'e
principal homog\`ene est plus restrictive que la
premi\`ere~\ref{XI.4.1}.
\end{remarques}

\subsection{}
\label{XI.4.7}
On obtient une notion encore plus restrictive de fibr\'e principal
homog\`ene en demandant que $S$ soit recouvert par des ouverts $S_i$
tels que pour tout~$i$, $P|S_i$ soit un fibr\'e \`a op\'erateurs
trivial sous~$G|S_i$: on dira alors que $P$ est un fibr\'e principal
homog\`ene \emph{localement trivial}.
\index{principal homog\`ene localement trivial (fibr\'e)|hyperpage}Les classes de ces fibr\'es, pour $G$ donn\'e, forment une partie
de~$\H^1(X,G)$, qui est en correspondance biunivoque avec
$\H^1(X,\cal{O}_X(G))$, o\`u $\cal{O}_X(G)$ est le faisceau (au sens
ordinaire) des sections de~$G$ sur~$S$, \cf \cite{XI.2}. Pour que ces $\H^1$
donnent encore lieu \`a une suite exacte de
cohomologie~\ref{XI.4.5}, il faut \'evidemment supposer que la suite
$0\to G'\to G\to G''\to 0$ soit exacte au sens raisonnable pour ce
nouveau contexte, \ie que $G$ soit fibr\'e localement trivial
sur~$G''$, de groupe $G'_{G''}$; cela signifie aussi que $u\colon
G'\to G$ est un noyau de~$v\colon G\to G''$, et que $G$ admet
localement une section sur~$G''$.

\subsection{}
\label{XI.4.8}
Il est \'evidemment tr\`es d\'esirable de continuer la suite
exacte~\ref{XI.4.5} en introduisant les groupes de cohomologie
sup\'erieurs $\H^i(X,G)$. Cela est possible en se pla\c cant dans
le cadre de la \og Cohomologie de Weil\fg: on consid\`ere la
cat\'egorie $\cal{B}$ des pr\'esch\'emas quasi-compacts sur~$S$,
muni de l'ensemble $\cal{M}$ des morphismes fid\`element plats et
quasi-compacts, qu'on appellera morphismes localisants. Un \og faisceau
de Weil\fg ab\'elien sur~$S$ (ou mieux, sur~$(\cal{B},\cal{M})$) est
alors un foncteur contravariant $\cal{F}$ de~$\cal{B}$ dans la
cat\'egorie des groupes ab\'eliens, transformant sommes en
produits, et une suite
${\def\labelstyle{\scriptstyle}\xymatrix@C=.5cm{{T''=T'\times_TT'}\ar@<+2pt>[rr]^-{\pr_1,\pr_2}\ar@<-2pt>[rr]&&{T'}\ar[r]^f&{T}}}$, avec $f\in\cal{M}$, en un diagramme \emph{exact}
d'ensembles $\xymatrix@C=.5cm{{\cal{F}(T)}\ar[r]&{\cal{F}(T')}\ar@<2pt>[r]\ar@<-2pt>[r]&{\cal{F}(T'')}}$. Les faisceaux de Weil
forment une cat\'egorie ab\'elienne \`a limites inductives
exactes admettant un g\'en\'erateur, donc admettant suffisamment
d'objets injectifs~\cite{XI.1}. Les foncteurs d\'eriv\'es droits du
foncteur $\Gamma(\cal{F})=\cal{F}(S)$ sont alors not\'es
$\H^i(S,\cal{F})$. D'autre part, tout $S$-groupe commutatif
d\'efinit \'evidemment un faisceau de Weil (VIII~\ref{VIII.5.2}),
dont le $\H^0$ et $\H^1$ ne sont autres que $\H^0(S,G)$ et
$\H^1(S,G)$, ce qui permet de d\'efinir les autres $\H^i(S,G)$ de
fa\c con raisonnable. On
montre d'ailleurs qu'une suite exacte de~$S$-groupes d\'efinit une
suite exacte de faisceaux de Weil, ce qui permet de retrouver et de
prolonger la suite exacte~\ref{XI.4.5}\footnote{Pour une \'etude
syst\'ematique de ce point de vue, \cf SGA~4~I~\`a~IX.}.

\subsection{}
\label{XI.4.9}
Il serait indiqu\'e de d\'evelopper les variantes non commutatives
de~\ref{XI.4.5} comme dans~\cite{XI.2}. Pour un d\'eveloppement
syst\'ematique, dans le cadre qui convient, des diverses notions
cohomologiques esquiss\'ees dans le pr\'esent num\'ero, nous
renvoyons \`a un travail en pr\'eparation de
J.\, \textsc{Giraud}\footnote{\Cf J.\, \textsc{Giraud}, 
\emph{Cohomologie non abélienne},
Springer-Verlag 1971.}.

\section{Cas particuliers de fibr\'es principaux}
\label{XI.5}

Supposons maintenant que $S$ soit connexe, et muni d'un point
g\'eom\'etrique~$a$, d'o\`u un groupe fondamental $\pi_1(S,a)$
permettant de classifier les rev\^etements \'etales de~$S$: la
cat\'egorie des rev\^etements \'etales de~$S$ est
\'equivalente \`a la cat\'egorie des ensembles finis o\`u
$\pi_1$ op\`ere contin\^ument. Il s'ensuit qu'un sch\'ema en
groupes fini et \'etale $G$ sur~$S$ est d\'etermin\'e
essentiellement par un groupe fini ordinaire $\cal{G}$, sur lequel
$\pi_1$ op\`ere contin\^ument par automorphismes de groupe. Un
rev\^etement \'etale $P$ de~$S$ o\`u $G$ op\`ere \`a droite
est d\'etermin\'e essentiellement par un ensemble fini $\cal{P}$
o\`u $\pi_1$ op\`ere contin\^ument (\`a gauche), et sur lequel
$\cal{G}$ op\`ere \`a droite de fa\c con compatible avec les
op\'erations de~$\pi_1$:
$$
s(p\cdot g)=(sp)\cdot(sg)\qquad \text{pour }s\in\pi_1,\
p\in\cal{P},\ g\in\cal{G}.
$$
On v\'erifie que $P$ est un fibr\'e principal homog\`ene au sens
de~\ref{XI.4.1} si et seulement si $\cal{P}$ est un ensemble principal
homog\`ene sous~$\cal{G}$ (utiliser par exemple le
crit\`ere~\ref{XI.4.2}). En d'autres termes, \emph{la catégorie
des fibrés principaux homogènes sur~$S$ de groupe $G$ est
équivalente à la catégorie des fibrés principaux
homogènes de groupe $G$ dans la catégorie des ensembles finis
où $\pi_1$ opère continûment.} On en d\'eduit en
particulier une bijection canonique, fonctorielle en $G$:
\begin{equation*}
\label{eq:XI.5.etoile}
\tag{$*$} \H^1(S,G)\isomto \H^1(\pi_1,\cal{G})\quoi,
\end{equation*}
o\`u
le deuxi\`eme membre d\'esigne l'ensemble des classes, \`a
isomorphisme pr\`es, des fibr\'es principaux homog\`enes sous
$\cal{G}$ dans la cat\'egorie des ensembles finis o\`u $\pi_1$
op\`ere (inutile d'ailleurs de pr\'eciser: contin\^ument),
ensemble qui s'explicite de fa\c con bien connue comme ensemble
quotient de l'ensemble $Z^1(\pi_1,G)$ des $1$-cocycles
$\varphi:\pi_1\to \cal{G}$ (satisfaisant $\varphi(1)=1$,
$\varphi(st)=\varphi(s)(s.\varphi(t))$) par le groupe $\cal{G}$ qui y
op\`ere de fa\c con naturelle).

\label{page-300}
Un cas important est celui o\`u $\pi_1$ op\`ere trivialement dans~$\cal{G}$, \ie lorsque $G$ est un rev\^etement compl\`etement
d\'ecompos\'e de~$S$, isomorphe \`a la somme de~$\cal{G}$
exemplaires de~$S$; on \'ecrit alors aussi $\H^1(S,\cal{G})$ au lieu
de~$\H^1(S,G)$, et cet ensemble est en correspondance biunivoque
\eqref{eq:XI.5.etoile} avec $\H^1(\pi_1,\cal G)= \Hom(\pi_1,\cal{G})/$automorphismes int\'erieurs de~$\cal{G}$. On notera d'ailleurs que dans ce cas, un fibr\'e
principal homog\`ene sur~$S$ de groupe $G$ n'est autre chose qu'un
\emph{revêtement principal} de~$S$ de groupe $\cal{G}$
(V~\ref{V.2.7}), et la correspondance biunivoque pr\'ec\'edente
est celle qui se d\'eduit de la correspondance entre rev\^etements
principaux de~$S$ de groupe~$\cal{G}$, \emph{ponctués} au-dessus
de~$a$, et les homomorphismes continus de~$\pi_1(S,a)$ dans $\cal{G}$
(\ref{V}~fin~du~\No \ref{V.5}).

L'int\'er\^et de relier la th\'eorie des rev\^etements
\'etales avec celle des fibr\'es principaux (d\'ej\`a
implicite dans \textsc{A.\, Weil}, G\'en\'eralisation des Fonctions
Ab\'eliennes, et explicit\'ee par \textsc{S.\, Lang} dans sa th\'eorie
g\'eom\'etrique du corps de classes, \cf Serre~\cite{XI.5}), vient du fait
que tout $S$-groupe qui est fini et \'etale sur~$S$ peut se plonger
dans un $S$-groupe~$H$, affine et lisse sur~$S$, \`a fibres
connexes, commutatif lorsque $G$ l'est de sorte que par la suite
exacte~\ref{XI.4.5} (et \'eventuellement ses variantes non
commutatives), la classification \og discr\`ete\fg des rev\^etements
principaux de groupe $G$ peut s'\'etudier \`a l'aide de la
classification \og continue\fg des fibr\'es principaux de groupe~$H$,
et r\'eciproquement d'ailleurs. Pour l'id\'ee de la construction
g\'en\'erale de l'immersion de~$G$ dans $H$ (assez peu
utilis\'ee en pratique semble-t-il), se reporter \`a
\cite[VI~2.8]{XI.5}. Nous nous contentons de d\'evelopper au \No suivant
deux cas particuliers importants, d'ailleurs classiques. Nous y aurons
besoin d'un r\'esultat auxiliaire:

\begin{proposition}
\label{XI.5.1}
Soit
$S$ un pr\'esch\'ema, $G$ un $S$-groupe isomorphe \`a
$\Gl(n)_S$ (par exemple $\GG_{m\,S}$) ou $\GG_{a\,S}$, alors tout
fibr\'e principal homog\`ene sous $G$ est localement trivial.
\end{proposition}

Pr\'ecisons que $\Gl(n)_S$ ($n$ entier $\geq 0$) d\'esigne le
$S$-groupe qui repr\'esente le foncteur contravariant $T\mto
\Gl(n,\Gamma(T,\cal{O}_T))$ en le $S$-pr\'esch\'ema~$T$, en
particulier $\GG_{m\,S}$ (\og groupe multiplicatif sur~$S$\fg)
repr\'esente le foncteur contravariant $T\mto
\Gamma(T,\cal{O}_T^*)$, donc comme pr\'esch\'ema sur~$S$ est
isomorphe \`a $\Spec {\cal{O}_S[t,t^{-1}]}$, o\`u $t$ est une
ind\'etermin\'ee. De m\^eme $\GG_{a\,S}$ repr\'esente le
foncteur contravariant $T\mto \Gamma(T,\cal{O}_T)$, il est donc
isomorphe comme $S$-pr\'esch\'ema \`a $\Spec(\cal{O}_S[t])$,
o\`u $t$ est une ind\'etermin\'ee. Notons que par d\'evissage,
\ref{XI.5.1} redonne le r\'esultat de locale trivialit\'e de
Rosenlicht, relatif au cas o\`u $G$ admet une \og suite de
composition\fg dont les facteurs cons\'ecutifs sont des groupes du
type envisag\'e ici. (Pour une \'etude plus fine des questions de
locale trivialit\'e des fibr\'es principaux homog\`enes,
\cf \cite{XI.7}~et~\cite{XI.3}).

La premi\`ere assertion se d\'emontre en remarquant que
$G(T)=\Aut(\cal{O}_T^n)$, et que les morphismes $S'\to S$ intervenant
dans~\ref{XI.4.1} (\ie qui sont fid\`element plats et
quasi-compacts) sont des morphismes de descente effective pour la
cat\'egorie fibr\'ee des Modules localement isomorphes \`a
$\cal{O}_T^n$, \ie localement libres de rang~$n$
(VIII~\ref{VIII.1.12}). La deuxi\`eme se d\'emontre de fa\c con
analogue, en notant que dans ce cas on a $G(T)=\Aut(\cal{E}_T)$,
o\`u $\cal{E}_T$ est l'\emph{extension} triviale de~$\cal{O}_T$ par
$\cal{O}_T$ (et o\`u les automorphismes bien entendu doivent
respecter la structure d'extension), et que les morphismes $S'\to S$
intervenant dans~\ref{XI.4.1} sont des morphismes de descente
effective pour la cat\'egorie fibr\'ee des extensions de~$\cal{O}_T$ par $\cal{O}_T$ (comme il r\'esulte facilement de
VIII~\ref{VIII.1.1}), et que de telles extensions sont automatiquement
localement triviales.

\begin{remarque}
\label{XI.5.2}
On notera que le m\^eme type de d\'emonstration s'applique au
groupe symplectique $\Symp(2n)_S$, compte tenu qu'une forme
altern\'ee sur un module localement isomorphe \`a
$\cal{O}_S^{2n}$, qui est \og non d\'eg\'en\'er\'ee\fg
\ie d\'efinit un isomorphisme de ce Module sur son dual, est
localement isomorphe \`a la forme standard. Le r\'esultat pour le
groupe orthogonal est par contre faux, d\'ej\`a si $S$ est le
spectre d'un corps, car il peut y avoir des formes
quadratiques sur un corps qui ne sont pas isomorphes \`a la forme
standard. D'ailleurs on montre essentiellement dans \cite{XI.3} que les
groupes $\Gl$, $\Symp$, $\GG_a$ et ceux qui se d\'evissent en tels
groupes, sont \`a peu de choses pr\`es les seuls pour lesquels on
ait un r\'esultat de trivialit\'e locale du type consid\'er\'e
ici.
\end{remarque}

\begin{corollaire}
\label{XI.5.3}
On a des bijections canoniques
$$
\H^1(S,\GL(n)_S)\isomfrom \H^1(S,\GL(n,\cal{O}_S))\quoi,
$$
en particulier
$$
\H^1(S,\GG_{m\,S})\isomfrom \H^1(S,\cal{O}_S^*)\quoi,
$$
et
$$
\H^1(S,\GG_{a\,S})\isomfrom \H^1(S,\cal{O}_S)\quoi,
$$
o\`u les deuxi\`emes membres d\'esignent des cohomologies de
l'espace topologique $S$ \`a coefficients dans des faisceaux
ordinaires.
\end{corollaire}

En particulier, $\H^1(S,\GL(n)_S)$ s'identifie \`a l'ensemble des
classes, \`a isomorphisme pr\`es, de Modules localement libres de
rang $n$ sur~$S$, et $\H^1(S,\GG_{a\,S})$ s'identifie \`a l'ensemble
des classes d'extensions du Module $\cal{O}_S$ par lui-m\^eme.

\section[Application aux rev\^etements principaux]{Application aux rev\^etements principaux: th\'eories de
Kummer et d'Artin-Schreier}
\label{XI.6}


\begin{proposition}
\label{XI.6.1}
Soient $S$ un pr\'esch\'ema, $n$ un entier $>0$, soit $u_n\colon
\GG_{m\,S}\to\GG_{m\,S}$ l'homomorphisme de puissance \nieme, et
$\bbmu_{n\,S}$
\label{indnot:kb}son noyau. Alors $\bbmu_{n\,S}$ est fini et localement libre de
rang $n$ sur~$S$, et il est \'etale sur~$S$ si et seulement si pour
tout $s\in S$, la caract\'eristique de~$s$ est premi\`ere \`a
$n$. La suite d'homomorphismes
$$
0\to \bbmu_{n\,S}\to\GG_{m\,S}\lto{u_n}\GG_{m\,S}\to 0
$$
est exacte au sens du \No \ref{XI.4}. (On l'appellera la \emph{suite exacte de
Kummer}
\index{Kummer (suite exacte de)|hyperpage}sur~$S$, relativement \`a l'entier~$n$).
\end{proposition}

On
a
$$
\GG_m=\Spec\cal{O}_S[t,t^{-1}],
$$
et $u_n$ correspond \`a l'homomorphisme $u_n$ sur les
$\cal{O}_S$-alg\`ebres affines, donn\'e par
$$
u_n(t)=t^n,
$$
d'autre part la section unit\'e de~$\GG_{m\,S}$ correspond \`a
l'homomorphisme d'augmentation de~$\cal{O}_S$-alg\`ebres, donn\'e
par
$$
\varepsilon(t)=1,
$$
dont le noyau est donc l'Id\'eal principal $(t-1)$. L'image de ce
dernier par $u_n$ est donc l'Id\'eal principal $(1-t^n)$, et on
trouve:
$$
\bbmu_{n\,S}=\Spec\cal{O}_S[t]/(1-t^n),
$$
ce qui montre en particulier que $\bbmu_{n\,S}$ est fini sur~$S$,
et d\'efini par une Alg\`ebre sur~$S$ qui est libre de rang $n$,
ayant la base form\'ee des $t^i$ ($0\leq i\leq n-1$). Pour qu'il soit
\'etale en $s\in S$, il faut et il suffit que l'Alg\`ebre
r\'eduite $k[t]/(1-t^n)$, o\`u $k=\kres(s)$, obtenue par adjonction
formelle des racines \niemes de l'unit\'e \`a $k$, soit
s\'eparable sur~$k$, \ie les racines de~$1-t^n$ dans une
cl\^oture alg\'ebrique de~$k$ sont toutes distinctes, ce qui
\'equivaut au fait que $n$ soit premier \`a la
caract\'eristique. Enfin, pour montrer que la suite
d'homomorphismes dans~\ref{XI.6.1} est exacte, on est ramen\'e en
vertu du crit\`ere~\ref{XI.4.2} \`a prouver que
$u_n$
est fid\`element plat. On peut \'evidemment supposer $S$ affine
d'anneau~$A$, donc $\GG_{m\,S}$ affine d'anneau $B=A[t,t^{-1}]$, et il
suffit de v\'erifier que $u_n$ fait de~$B$ un module libre de rang
$n$ sur~$B$, ou ce qui revient au m\^eme, que $u_n$ est injectif, et
que $A[t,t^{-1}]$ est un module libre de rang $n$ sur~$A[t^n,t^{-n}]$. En effet, on v\'erifie facilement que les
$t^i$ ($0\leq i\leq n-1$) forment une base de l'un sur l'autre, ce qui
ach\`eve la d\'emonstration.

\begin{definition}
\label{XI.6.2}
On
appelle $\bbmu_{n\,S}$ le \emph{groupe de Kummer de rang $n$
sur~$S$}, et on appelle \emph{revêtement principal kummérien
de rang
$n$ sur~$S$}
tout fibr\'e principal homog\`ene sur~$S$ de groupe le groupe de
Kummer de rang~$n$.
\end{definition}

L'ensemble de ces rev\^etements est un groupe, not\'e
$\H^1(S,\bbmu_{n\,S})$ ou simplement $\H^1(S,\bbmu_n)$. On
notera que la formation du groupe de Kummer de rang $n$ sur~$S$ est
compatible avec l'extension de la base, donc que $\bbmu_{n\,S}$
provient par extension de la base du \emph{groupe de Kummer absolu
$\bbmu_n$} sur~$\Spec(\ZZ)$.

D\'esignons par $(\ZZ/n\ZZ)_S$ le $S$-groupe d\'efini par le
groupe fini ordinaire $\ZZ/n\ZZ$. Si~$G$ est un $S$-groupe quelconque,
les homomorphismes de~$S$-groupes $u$ de~$(\ZZ/n\ZZ)_S$ dans $G$
correspondent biunivoquement, et de fa\c con compatible avec le
changement de base, aux sections de~$G$ sur~$S$ dont la puissance
\nieme est la section unit\'e, en faisant correspondre \`a
$u$ l'image par $u$ de la section de~$(\ZZ/n\ZZ)_S$ sur~$S$ d\'efini
par le g\'en\'erateur $1\bmod n\ZZ$ de~$\ZZ/n\ZZ$. Ceci pos\'e:

\begin{corollaire}
\label{XI.6.3} Si $\bbmu_{n\,S}$ est \'etale sur~$S$, on
obtient ainsi une correspondance biunivoque entre les isomorphismes de~$S$-groupes $(\ZZ/n\ZZ)_S\isomto\bbmu_{n\,S}$, et les sections de~$\cal{O}_S$ qui sont d'ordre $n$ exactement sur chaque composante
connexe de~$S$ (une telle section s'appellera \og racine primitive
\nieme de l'unit\'e sur~$S$\fg). Donc pour que
$\bbmu_{n\,S}$ soit isomorphe en tant que $S$-groupe \`a
$(\ZZ/n\ZZ)_S$, il faut et il suffit qu'il soit \'etale
sur~$S$,
\ie que les caract\'eristiques r\'esiduelles de~$S$ soient
premi\`eres \`a $n$, et qu'il existe une racine primitive
\nieme de l'unit\'e sur~$S$.
\end{corollaire}

Cela explique le r\^ole jou\'e dans la th\'eorie kumm\'erienne
classique par l'hypoth\`ese que le corps de base (jouant le r\^ole
de~$S$) soit de caract\'eristique premi\`ere \`a $n$ et
contienne les racines \niemes de l'unit\'e, et par le choix
d'une racine primitive \nieme de l'unit\'e. Une fois qu'on
dispose du langage des sch\'emas, il n'y a plus lieu de s'embarrasser
de ces hypoth\`eses, et il convient de raisonner directement sur~$\bbmu_n$ au lieu de~$\ZZ/n\ZZ$. Ainsi, la conjonction
de~\ref{XI.6.1}, \ref{XI.4.5} et~\ref{XI.5.3} nous donne la relation
g\'en\'erale suivante entre la th\'eorie des rev\^etements
principaux kumm\'eriens et celle des groupes de Picard:

\begin{proposition}
\label{XI.6.4}
Soient
$S$ un pr\'esch\'ema, on a une suite exacte canonique
$$
\textstyle
0\to\H^0(S,\bbmu_n)\to\H^0(S,\cal{O}^*_S)\to
\H^0(S,\cal{O}^*_S)\to\H^1(S,\bbmu_n)\to
\H^1(S,\cal{O}^*_S)\to\H^1(S,\cal{O}^*_S),
$$
d'o\`u, en posant
$\H^1(S,\cal{O}_S^*)=\Pic(S)$,
et en d\'esignant pour tout groupe ab\'elien $A$, par ${_n}A$ et
$A_n$ les noyau et conoyau de la multiplication par $n$ dans $A$, la
suite exacte:
$$
0\to\H^0(S,\cal{O}_S)^*_n\to\H^1(S,\bbmu_n)\to
{{}_n\Pic}(S)\to 0.
$$
\end{proposition}

Nous allons expliciter deux cas importants, o\`u l'un ou l'autre
terme extr\^eme de cette suite exacte sont nuls:

\begin{corollaire}
\label{XI.6.5} Supposons ${}_n\Pic(S)=0$, (par exemple
que $S$ soit le spectre d'un anneau local, ou d'un anneau factoriel),
et soit $A$ l'anneau $\H^0(S,\cal{O}_S)$. Alors on a un isomorphisme
canonique
$$
\H^1(S,\bbmu_n)\isomto A^*/{A^*}^n.
$$
\end{corollaire}

C'est essentiellement l'\'enonc\'e classique de la th\'eorie de
Kummer, lorsque $S$ est le spectre d'un corps.

\begin{corollaire}
\label{XI.6.6} Supposons que tout \'el\'ement de
$\H^0(S,\cal{O}_S)$ soit une puissance \nieme, par exemple que
$\H^0(S,\cal{O}_S)$ soit un compos\'e de corps alg\'ebriquement
clos ou que $S$ soit r\'eduit et propre sur un corps
alg\'ebriquement clos $k$. Alors on a un isomorphisme canonique
$$
\H^1(S,\bbmu_n)\isomto{}_n\Pic(S).
$$
\end{corollaire}

En particulier, lorsque $S$ est propre et connexe sur un corps
alg\'ebriquement clos $k$, cela met en relation le groupe
fondamental de~$S$ avec les points d'ordre fini du sch\'ema de
Picard $P$ de~$S$ sur~$k$; ainsi on aura un isomorphisme
$$
\Hom(\pi_1(S),\ZZ/n\ZZ)\simeq {_n P}(k)
$$
pour $n$ premier \`a la caract\'eristique, qui est souvent
utilis\'e en g\'eom\'etrie
alg\'ebrique. Comme application, lorsque la composante connexe $P^0$
de~$P$ est un sch\'ema en groupes complet, de dimension $g$, on voit
en utilisant les r\'esultats rappel\'es dans le \No \ref{XI.2}, et la
finitude du groupe de torsion de
N\'eron-Severi,
que pour
tout nombre premier $\ell$ premier \`a la caract\'eristique, la
composante $\ell$-primaire du groupe fondamental $\pi_1(S)$ rendu
ab\'elien est un module de type fini et de rang $2g$ sur l'anneau
$\ZZ_\ell$ des entiers $\ell$-adiques (et d'ailleurs libre sauf pour
un nombre fini au plus de valeurs de~$\ell$). Comme l'a remarqu\'e
Serre, cela permet de prouver sous certaines conditions que lorsque
$X$ est un sch\'ema plat et projectif sur~$S$ connexe, alors les
sch\'emas de Picard des fibres de~$X$ ont toutes la m\^eme
dimension, en appliquant le th\'eor\`eme de semicontinuit\'e
(X~\ref{X.2.3}); l'argument de Serre s'applique d\`es que le
sch\'ema de Picard de~$X$ sur~$S$ existe, et que les Picards
connexes des fibres de~$X$ sur~$S$ sont des sch\'emas en groupes
propres, par exemple lorsque les fibres g\'eom\'etriques de~$X$
sur~$S$ sont normales ($X$ \'etant toujours plat et projectif
sur~$S$), en particulier si $X$ est lisse et projectif sur~$S$.

Soit maintenant $p$ un nombre premier, et supposons que $S$ soit un
pr\'esch\'ema de caract\'eristique $p$, \ie tel que
$p\cdot\cal{O}_S=0$. Alors l'homomorphisme de puissance \pieme
dans $\cal{O}_S$ est additif, et le morphisme correspondant, obtenu en
rempla\c cant $S$ par un~$T$ variable sur~$S$:
$$
F\colon\GG_{a\,S}\to\GG_{a\,S}
$$
est donc un homomorphisme de~$S$-groupes, qu'on appelle
l'\emph{homomorphisme de Frobenius} (N.B. Un tel morphisme est
d\'efini pour tout $S$-pr\'esch\'ema $G$ qui provient par
extension de la base d'un pr\'esch\'ema $G_0$ sur le corps premier
$\ZZ/p\ZZ$, et ce morphisme est un homomorphisme de groupes si $G_0$
est un pr\'esch\'ema en groupes). Nous poserons:
$$
\wp=\id-F\colon\GG_{a\,S}\to\GG_{a\,S}.
$$
Consid\'erons d'autre part le $S$-groupe $(\ZZ/p\ZZ)_S$ d\'efini
par le groupe fini ordinaire~$\ZZ/p\ZZ$, nous avons dit que pour tout
$S$-groupe~$G$, les homomorphismes de~$S$-groupe de~$(\ZZ/p\ZZ)_S$
dans $G$ correspondent biunivoquement aux sections de~$G$ sur~$S$ dont
la puissance \pieme est la section unit\'e. Lorsque
$G=\GG_{a\,S}$, elles correspondent
donc aux sections quelconques de~$G$ sur~$S$. Prenant en particulier
la section de~$\GG_{a\,S}$ sur~$S$ correspondant \`a la section
unit\'e du faisceau d'anneaux $\cal{O}_S$, on trouve un
homomorphisme de~$S$-groupes
$$
i\colon(\ZZ/p\ZZ)_S\to\GG_{a\,S}
$$

\begin{proposition}
\label{XI.6.7} La suite d'homomorphismes de~$S$-groupes
$$
0\to(\ZZ/p\ZZ)_S\to\GG_{a\,S}\to \GG_{a\,S}\to 0
$$
est exacte (au sens du \No \ref{XI.4}). (On l'appelle la suite exacte
\emph{d'Artin-Schreier} sur~$S$).
\index{Artin-Schreier (suite exacte d')|hyperpage}\end{proposition}

Il suffit de le prouver sur le corps premier $k=\ZZ/p\ZZ$. Il suffit
de remarquer que l'homomorphisme $\wp^*\colon k[t]\to k[t]$ d\'efini
par $\wp^*(t)=t-t^p$ fait de~$k[t]$ un module libre de rang $p$ sur~$k[t]$, de fa\c con pr\'ecise que $k[t]$ est un module libre sur~$k[s]$, o\`u $s=t-t^p$, ayant la base form\'ee des $t^i$ ($0\leq i\leq p-1$).

On en conclut, utilisant \ref{XI.4.5} et~\ref{XI.5.3}:

\begin{proposition}
\label{XI.6.8} On a une suite exacte canonique:
\begin{multline*}
0\to\H^0(S,\ZZ/p\ZZ)\to\H^0(S,\cal{O}_S)\to\H^0(S,\cal{O}_S)\to
\H^1(S,\ZZ/p\ZZ)\\
\to\H^1(S,\cal{O}_S)\to\H^1(S,\cal{O}_S),
\end{multline*}
d'o\`u une suite exacte:
$$
0\to\H^0(S,\cal{O}_S)/\wp\H^0(S,\cal{O}_S)\to
\H^1(S,\ZZ/p\ZZ)\to\H^1(S,\cal{O}_S)^F\to 0,
$$
(o\`u l'exposant $F$ dans le dernier terme signifie le sous-groupe
des invariants par l'endomorphisme $F$, \'egal au noyau de
$\wp=\id-F$).
\end{proposition}

Explicitons encore deux cas extr\^emes:

\begin{corollaire}
\label{XI.6.9} Supposons que $\H^1(S,\cal{O}_S)^F=0$, par exemple
que $S$ soit un sch\'ema affine. Alors, posant
$A=\H^0(S,\cal{O}_S)$, on a un isomorphisme canonique
$$
\H^1(S,\ZZ/p\ZZ)\isomto A/\wp A.
$$
\end{corollaire}

C'est la \emph{théorie d'Artin-Schreier} dans la forme classique,
du moins lorsque $A$ est le spectre d'un corps.

\begin{corollaire}
\label{XI.6.10}
Supposons
que $\wp\H^0(S,\cal{O}_S){=}\H^0(S,\cal{O}_S)$, par exemple
que $\H^0(S,\cal{O}_S)$ soit un compos\'e de corps
alg\'ebriquement clos, ou que $S$ soit propre sur un corps
alg\'ebriquement clos. Alors on a un isomorphisme canonique:
$$
\H^1(S,\ZZ/p\ZZ)\isomto\H^1(S,\cal{O}_S)^F
$$
\end{corollaire}

\begin{remarques}
\label{XI.6.11}
Le dernier \'enonc\'e est d\^u \`a
J-P.\, Serre
\cite{XI.9}. Il est
possible \'egalement de d\'evelopper une th\'eorie analogue pour
le groupe structural $\ZZ/p^n\ZZ$ pour $n$ quelconque, en utilisant au
lieu de~$\GG_a$ le sch\'ema en groupes de Witt $\WW_n$, \cf \loccit On notera qu'en caract\'eristique $p>0$, la th\'eorie de
Kummer ne donne plus de renseignement sur les rev\^etements
principaux d'ordre $p$, puisque $\bbmu_p$ est alors un groupe
\og infinit\'esimal\fg \ie radiciel sur la base, donc sans rapport
direct avec $\ZZ/p\ZZ$; aussi \`a premi\`ere vue, la th\'eorie
de ces rev\^etements n'est plus justiciable (lorsque $S$ est un
sch\'ema propre sur un corps alg\'ebriquement clos pour fixer les
id\'ees), de la th\'eorie du sch\'ema de Picard comme
dans~\ref{XI.6.6}. N\'eanmoins, si on se rappelle que l'espace
tangent de Zariski \`a l'origine dans
$\mathbf{Pic}_{S/K}$\footnote{Pour la d\'efinition de
$\mathbf{Pic}_{S/K}$, \cf A. Grothendieck, S\'em. Bourbaki \No 232
(F\'evrier 1962).} s'identifie \`a $\H^1(S,\cal{O}_S)$, on
constate que \emph{la connaissance du schéma en groupes
${_p}\mathbf{Pic}_{S/k}$, noyau de la multiplication par $p$ dans
$\mathbf{Pic}_{S/k}$, implique celle de~$\H^1(S,\ZZ/p\ZZ)$ aussi bien
que celle de~$\H^1(S,\bbmu_p)$; on notera qu'elle implique aussi
celle de~$\H^1(S,\bbalpha_p)$}, o\`u $\bbalpha_p$
d\'esigne le sch\'ema en groupes infinit\'esimal sur le corps
premier, noyau de~$F\colon\GG_a\to\GG_a$ (qui peut se d\'ecrire
aussi comme le spectre de l'alg\`ebre enveloppante restreinte de la
$p$-alg\`ebre de Lie triviale de dimension~$1$): en effet la suite
exacte~\ref{XI.4.5} donne ici:
$$
\H^1(S,\bbalpha_p)\simeq\Ker
(F\colon\H^1(S,\cal{O}_S)\to\H^1(S,\cal{O}_S)),
$$
et plus g\'en\'eralement, d\'esignant par $\bbalpha_{p^n}$
le noyau dans $\GG_a$ du \nieme it\'er\'e de~$F$, on aura
$$
\H^1(S,\bbalpha_{p^n})\simeq\Ker(F^n\colon\H^1(S,\cal{O}_S)\to
\H^1(S,\cal{O}_S)).
$$

En fait, la connaissance de~${_p}\mathbf{Pic}_{S/k}$ \'equivaut \`a
celle de~$\H^1(S,G)$ pour tout groupe alg\'ebrique
commutatif fini annul\'e par $p$, plus g\'en\'eralement, la
connaissance de~${_{p^n}}\mathbf{Pic}_{S/k}$ \'equivaut \`a celle
de~$\H^1(S,G)$ pour tout groupe alg\'ebrique commutatif fini $G$
annul\'e par $p^n$, en vertu du th\'eor\`eme suivant qui englobe
dans le cas envisag\'e \`a la fois la th\'eorie de Kummer et
celle de Artin-Schreier:

Soit $G$ un groupe alg\'ebrique fini sur~$k$,
$D(G)=\SheafHom_{k\textup{-groupes}}(G,\GG_m)$ son \emph{dual de
Cartier} (dont l'alg\`ebre affine est port\'ee par l'espace
vectoriel dual de l'alg\`ebre affine de~$G$, \ie par
l'hyperalg\`ebre de~$G$ au sens de Dieudonn\'e-Cartier), alors on
a un isomorphisme canonique:
\begin{equation}\label{eqXI.6.11}\tag{$*$}
\H^1(S,G)\simeq\Hom_{k\textup{-groupes}}(D(G),\mathbf{Pic}_{S/k}).
\end{equation}
(N.B. $S$ est un sch\'ema propre sur~$k$ alg\'ebriquement clos, tel
que $\H^0(S,\cal{O}_S)=k$). Cette formule peut encore s'exprimer en
disant que le \og vrai groupe fondamental\fg de~$S$ auquel il \'etait
fait allusion au \No \ref{XI.2}, rendu ab\'elien, est isomorphe \`a la
limite projective des $D(P_i)$, o\`u $P_i$ parcourt les sous-groupes
alg\'ebriques \emph{finis} de~$\mathbf{Pic}_{S/k}$, qu'on notera
$T^\bbullet(\mathbf{Pic}_{X/k})$. Lorsque $S$ est une vari\'et\'e
ab\'elienne, on a vu dans~\ref{XI.2.1} que ce groupe est
\'egalement isomorphe au \og vrai\fg module de Tate
$T_{\bbullet}(S)=\varprojlim\, {{_n}S}$, et l'isomorphisme
\eqref{eqXI.6.11} s'\'ecrit alors de fa\c con plus frappante
$$
\Ext^1(A,G)\simeq\Hom(D(G),B),
$$
$A$ \'etant une vari\'et\'e ab\'elienne, $B$ sa duale, $G$ un
groupe alg\'ebrique fini sur~$k$. Les r\'esultats qu'on vient
d'indiquer peuvent se g\'en\'eraliser d'ailleurs au cas o\`u $k$
est remplac\'e par un pr\'esch\'ema de base quelconque, et \`a
d'autres groupes de coefficients $G$ que des groupes finis.
\end{remarques}

\begin{thebibliography}{10}{XI.7}

\bibitem[1]{XI.1} A.\, \textsc{Grothendieck},
Sur quelques points d'alg\`ebre homologique,
T\^ohoku Math.~J.\ \textbf{9} (1957) p.\, 119--221.

\bibitem[2]{XI.2} A.\, \textsc{Grothendieck}, A general theory of fibre spaces with
structure sheaf, University of Kansas, 1955.

\bibitem[3]{XI.3} A.\, \textsc{Grothendieck}, Torsion homologique et sections
rationnelles, S\'eminaire \hbox{Chevalley}, 16 juin 1958.

\bibitem[4]{XI.4} S.\, \textsc{Lang}, Abelian varieties, Interscience tracts in pure
and applied mathematics, \No 7, New York.

\bibitem[5]{XI.5}
J-P.\,\textsc{Serre}, Groupes alg\'ebriques et corps de classes,
Actualit\'es Scient. et Ind.\ \No 1264, Hermann, Paris, 1959.

\bibitem[6]{XI.6}
J-P.\,\textsc{Serre}, Groupes proalg\'ebriques,
Publications Math\'ematiques de l'IHES \textbf{7} (1960), p.\, 1--67.

\bibitem[7]{XI.7}
J-P.\,\textsc{Serre}, Espaces fibr\'es alg\'ebriques,
S\'eminaire Chevalley, 21 avril 1958.

\bibitem[8]{XI.8}
J-P.\,\textsc{Serre}, Quelques propri\'et\'es des
vari\'et\'es ab\'eliennes en
caract\'eristique~$p$, Amer.\ J.\ Math. \textbf{80} (1958), p.\, 715--739.

\bibitem[9]{XI.9}
J-P.\,\textsc{Serre}, Sur la topologie des vari\'et\'es
alg\'ebriques en
caract\'eristique~$p$,
Symposium Internacional de Topologia
Algebraica, 1958.

\bibitem[10]{XI.10}
J-P.\,\textsc{Serre}, On the fundamental group of a unirational
variety,
J.~London Math.\ Soc.\ \textbf{34} (1959), p.\, 481--484.
\end{thebibliography}


\chapterspace{-4}
